%! AP Statistics — Unit 4, Lesson 4.12 — Group Activity
\documentclass[10pt]{article}
\usepackage{apstats-article}
\usepackage{apstats-boxes}

\begin{document}

\pageheader{Unit 4, Lesson 4.12}{Group Activity: Geometric vs.\ Binomial}
\namepartnerperiod

\vspace{0.1in}

\begin{scenariobox}[Directions]{sky}
\small
Work through all three tiers. Show all formula setup. Use complete sentences
when interpreting parameters.
\end{scenariobox}

\vspace{0.1in}

% ── TIER R ────────────────────────────────────────────────────────────────────
\begin{tcolorbox}[colback=goldbg, colframe=goldacc, boxrule=1pt, arc=2mm,
    title={\bfseries Tier R \quad Remediate: Identifying the Distribution},
    fonttitle=\small]
\small

For each scenario, (a) identify the distribution (binomial, geometric, or neither),
(b) state the key distinction, and (c) if geometric or binomial, state $X$ and the
parameters.

\vspace{10pt}
\textbf{Scenario 1:} A coin is flipped until tails appears.
$X =$ the number of flips.
\begin{enumerate}[label=(\alph*), leftmargin=2em, itemsep=6pt]
  \item Distribution: \blank{4cm}
  \item Key distinction: \writeline
  \item Parameters: \blank{6cm}
\end{enumerate}

\vspace{10pt}
\textbf{Scenario 2:} A coin is flipped 10 times.
$X =$ the number of tails.
\begin{enumerate}[label=(\alph*), leftmargin=2em, itemsep=6pt]
  \item Distribution: \blank{4cm}
  \item Key distinction: \writeline
  \item Parameters: \blank{6cm}
\end{enumerate}

\vspace{10pt}
\textbf{Scenario 3:} A tetrahedral die (4 sides) is rolled. The die is fair.
$X =$ the number of the first roll.
\begin{enumerate}[label=(\alph*), leftmargin=2em, itemsep=6pt]
  \item Distribution: \blank{4cm}
  \item Key distinction: \writeline
\end{enumerate}
\vspace{4pt}
\end{tcolorbox}

\vspace{0.1in}

% ── TIER A ────────────────────────────────────────────────────────────────────
\begin{tcolorbox}[colback=sky, colframe=navy, boxrule=1pt, arc=2mm,
    title={\bfseries Tier A \quad Approaching Proficiency: Full Geometric Analysis},
    fonttitle=\small]
\small

\textbf{Context:} An archer hits the bullseye 30\% of the time. She shoots arrows
one at a time until hitting the bullseye.
Let $X =$ the number of arrows shot until (and including) the first bullseye.

\vspace{8pt}
(a) Check BITS. Write one sentence per condition.

\vspace{4pt}
B: \writeline

I: \writeline

T: \writeline

S: \writeline

\vspace{6pt}
State the distribution: $X \sim $ \blank{4cm}

\vspace{8pt}
(b) Compute $P(X = 5)$. Show the formula and the numerical answer.
\[
  P(X = 5) = (\quad)^{4}(\quad) = \hspace{5cm} \approx \hspace{2cm}
\]

\vspace{8pt}
(c) Compute $\mu_X$ and interpret in context.
\vspace{4pt}
$\mu_X = $ \blank{4cm} $= $ \blank{2cm}

\vspace{4pt}
Interpretation: \writeline\\[1pt]\writeline

\vspace{8pt}
(d) Find $P(X \le 4)$ using \texttt{geometcdf}. What does this represent?
\vspace{4pt}
$P(X \le 4) = \texttt{geometcdf(}\blank{1cm}\texttt{,}\;4\texttt{)} \approx $ \blank{2cm}

\vspace{4pt}
Interpretation: \writeline\\[1pt]\writeline
\vspace{4pt}
\end{tcolorbox}

\vspace{0.1in}

% ── TIER E ────────────────────────────────────────────────────────────────────
\begin{tcolorbox}[colback=greenbg, colframe=greenacc, boxrule=1pt, arc=2mm,
    title={\bfseries Tier E \quad Extension: Distribution Table and Shape},
    fonttitle=\small]
\small

\textbf{Context (same archer):} $X \sim \text{Geom}(0.30)$.

\vspace{8pt}
(a) Build the partial distribution table for $X = 1, 2, 3, 4, 5$.

\vspace{4pt}
\renewcommand{\arraystretch}{1.8}
\begin{tabularx}{\linewidth}{@{} >{\bfseries}p{1.2cm} X X X X X @{}}
$k$ & 1 & 2 & 3 & 4 & 5 \\
\hline
$P(X=k)$ & \blank{2cm} & \blank{2cm} & \blank{2cm} & \blank{2cm} & \blank{2cm} \\
\end{tabularx}

\vspace{8pt}
(b) Describe the shape of the geometric distribution based on your table.
Does it look like the binomial distribution for the same $p$?
\vspace{4pt}
\writeline\\[1pt]\writeline

\vspace{8pt}
(c) Now consider $Y \sim B(5, 0.30)$ (same $p$, fixed $n = 5$).
Compare $P(X = 5)$ and $P(Y = 5)$. Explain why they are different,
and state what each probability means in context.
\vspace{4pt}
\writeline\\[1pt]\writeline\\[1pt]\writeline
\vspace{4pt}
\end{tcolorbox}

\end{document}
